\chapter{Θεωρητικό υπόβαθρο}
\label{ch:background}

Το κεφάλαιο συγκεντρώνει τις έννοιες στις οποίες στηρίζονται τα επόμενα: τι είναι η μικτή πραγματικότητα και τι προσφέρει η συσκευή που χρησιμοποιήθηκε, πώς περιγράφονται οι θέσεις των αισθητήρων, πώς ευθυγραμμίζονται δύο συστήματα συντεταγμένων, πώς εκτιμάται ένα μέγεθος ανάμεσα σε μετρήσεις, ποιοι περιορισμοί ισχύουν στην απόδοση γραφικών σε αυτόνομες συσκευές, και πώς οργανώνεται η παρακολούθηση κατασκευών.

\section{Μικτή πραγματικότητα}
\label{sec:bg-mr}

Οι \en{Milgram} και \en{Kishino} περιέγραψαν το \emph{συνεχές πραγματικότητας--εικονικότητας}, που εκτείνεται από το εντελώς πραγματικό περιβάλλον ως το εντελώς εικονικό \cite{milgram1994taxonomy}. Ενδιάμεσα βρίσκεται η \emph{μικτή πραγματικότητα} (MR), όπου πραγματικά και εικονικά στοιχεία συνυπάρχουν. Η επαυξημένη πραγματικότητα (AR) είναι η περιοχή όπου κυριαρχεί το πραγματικό περιβάλλον και τα εικονικά στοιχεία προστίθενται πάνω του.

Ο \en{Azuma} όρισε ένα σύστημα AR με τρία χαρακτηριστικά: συνδυάζει πραγματικό και εικονικό, είναι διαδραστικό σε πραγματικό χρόνο, και τα εικονικά στοιχεία είναι \emph{ευθυγραμμισμένα} (\en{registered}) στον τρισδιάστατο χώρο \cite{azuma1997survey}. Το τρίτο χαρακτηριστικό είναι και το δυσκολότερο: ένα εικονικό σημείο που «κολυμπά» μερικές δεκάδες εκατοστά μακριά από τον αισθητήρα που αντιπροσωπεύει μπορεί να παραπλανήσει περισσότερο από ό,τι βοηθά. Για τον λόγο αυτό η ευθυγράμμιση είναι κεντρικό θέμα της παρούσας εργασίας. Αναλυτική επισκόπηση του πεδίου δίνουν οι \en{Billinghurst} κ.ά. \cite{billinghurst2015survey}.

Οι συσκευές κεφαλής χωρίζονται σε δύο κατηγορίες. Στις συσκευές \emph{οπτικής διαφάνειας} (\en{optical see-through}, π.χ. HoloLens) ο χρήστης βλέπει τον πραγματικό κόσμο μέσα από ημιδιαφανή οπτικά στοιχεία, πάνω στα οποία προβάλλονται τα εικονικά. Στις συσκευές \emph{video passthrough} (π.χ. Quest~3) κάμερες καταγράφουν τον περιβάλλοντα χώρο και η εικόνα του αναπαράγεται στις οθόνες μαζί με τα εικονικά στοιχεία. Οι δεύτερες έχουν μεγαλύτερο οπτικό πεδίο και πλήρη έλεγχο της σύνθεσης (π.χ. μπορούν να σκοτεινιάσουν τον πραγματικό χώρο), με κόστος την ανάλυση και την καθυστέρηση της κάμερας.

\subsection{Η συσκευή Meta Quest 3}
\label{sec:bg-quest}

Το Meta Quest~3 είναι αυτόνομο headset με επεξεργαστή Qualcomm Snapdragon XR2 Gen~2 και λειτουργικό βασισμένο στο Android (Horizon OS). Για την παρούσα εργασία έχουν σημασία τα εξής χαρακτηριστικά:
\begin{itemize}
  \item \textbf{Εντοπισμός θέσης inside-out}: η συσκευή εκτιμά τη θέση και τον προσανατολισμό της με έξι βαθμούς ελευθερίας από τις κάμερες και τους αδρανειακούς αισθητήρες της (SLAM), χωρίς εξωτερικούς σταθμούς. Η κατακόρυφος είναι γνωστή με ακρίβεια από το επιταχυνσιόμετρο.
  \item \textbf{Έγχρωμο passthrough} από δύο κάμερες RGB, που επιτρέπει εφαρμογές μικτής πραγματικότητας.
  \item \textbf{Αισθητήρας βάθους}, μέσω του οποίου η εφαρμογή γνωρίζει την απόσταση των πραγματικών επιφανειών και μπορεί να «κρύψει» εικονικά αντικείμενα πίσω τους (απόκρυψη, \en{occlusion}).
  \item \textbf{Χειριστήρια} με εντοπισμό έξι βαθμών ελευθερίας, που χρησιμοποιούνται εδώ και ως όργανο σκόπευσης: η άκρη τους ακουμπά τα σημεία ελέγχου.
  \item \textbf{Ρυθμός ανανέωσης} 72, 90 ή 120\,Hz, και \textbf{foveated rendering} σταθερού τύπου.
\end{itemize}

\review{Επιβεβαίωσε τα τεχνικά χαρακτηριστικά του Quest~3 από \textbf{επίσημη πηγή} (σελίδα προϊόντος ή τεκμηρίωση της Meta) και πρόσθεσε αναφορά (\code{@online} στο \code{references.bib} με ημερομηνία πρόσβασης). Αν ο επιβλέπων θέλει πίνακα χαρακτηριστικών (ανάλυση, οπτικό πεδίο, μνήμη), πρόσθεσέ τον εδώ.}

\subsection{Spatial anchors}
\label{sec:bg-anchors}

Ο εντοπισμός θέσης της συσκευής ορίζει ένα σύστημα συντεταγμένων που ισχύει μόνο για την τρέχουσα συνεδρία: η αρχή του μπορεί να αλλάξει σε κάθε εκκίνηση ή επανακέντρωση. Για να «κολλήσει» ένα εικονικό στοιχείο σε ένα σημείο του πραγματικού χώρου με διάρκεια, οι πλατφόρμες προσφέρουν \emph{spatial anchors}: σημεία που η συσκευή συνδέει με τον χάρτη οπτικών χαρακτηριστικών του περιβάλλοντος. Ένα anchor αποθηκεύεται με ένα αναγνωριστικό (UUID) και, σε επόμενη συνεδρία, η συσκευή προσπαθεί να το \emph{επανεντοπίσει} (\en{relocalize}) αναγνωρίζοντας τον χώρο. Η επιτυχία και η ακρίβεια του επανεντοπισμού εξαρτώνται από τον φωτισμό, την υφή των επιφανειών και τις αλλαγές του χώρου από την προηγούμενη επίσκεψη. Επιπλέον, η ακρίβεια ενός anchor υποβαθμίζεται με την απόσταση από αυτό, κάτι που σε ένα μεγάλο εργοτάξιο οδηγεί στη χρήση περισσότερων του ενός anchors. Η συμπεριφορά αυτή μετριέται πειραματικά στην \cref{sec:exp-anchors}.

\section{Συστήματα συντεταγμένων}
\label{sec:bg-crs}

Οι θέσεις των αισθητήρων ενός έργου δίνονται από τοπογραφική αποτύπωση σε ένα \emph{προβολικό} σύστημα συντεταγμένων: ένα επίπεδο, μετρικό σύστημα με άξονες Ανατολή ($E$) και Βορράς ($N$), και ένα υψόμετρο ($H$). Στην Ελλάδα το συνηθέστερο είναι το ΕΓΣΑ87, ενώ σε άλλα έργα χρησιμοποιούνται ζώνες UTM ή τοπικοί κάναβοι του έργου. Το σύστημα που περιγράφεται εδώ δεν δένεται σε συγκεκριμένο σύστημα: δέχεται οποιονδήποτε μετρικό κάναβο $E/N/H$, αρκεί αισθητήρες, σημεία ελέγχου και μοντέλα να εκφράζονται στον ίδιο.

Οι μηχανές απόδοσης γραφικών χρησιμοποιούν δικό τους σύστημα. Η Unity έχει αριστερόστροφο σύστημα με τον άξονα $y$ προς τα πάνω, ενώ ο κάναβος του έργου είναι δεξιόστροφος με το $H$ προς τα πάνω. Η αντιστοίχιση από τον κάναβο $(E, N, H)$ στο τοπικό σύστημα $(x, y, z)$ είναι
\begin{equation}
  \begin{pmatrix} x \\ y \\ z \end{pmatrix} =
  \begin{pmatrix} E - E_0 \\ H - H_0 \\ N - N_0 \end{pmatrix},
  \label{eq:enh-to-local}
\end{equation}
όπου $(E_0, N_0, H_0)$ η αρχή του έργου. Η αφαίρεση της αρχής είναι απαραίτητη: οι συντεταγμένες ΕΓΣΑ87 έχουν τάξη μεγέθους $10^5$--$10^6$\,m, και σε αριθμούς κινητής υποδιαστολής απλής ακρίβειας (24 bit mantissa) αυτό αφήνει ακρίβεια λίγων εκατοστών. Με την αρχή μέσα στο έργο, οι τιμές μένουν μικρές και η ακρίβεια φτάνει το χιλιοστό. Η αρχή κρατιέται σε διπλή ακρίβεια.

\section{Χωρική ευθυγράμμιση}
\label{sec:bg-registration}

\emph{Ευθυγράμμιση} (\en{registration}) είναι η εκτίμηση του μετασχηματισμού που φέρνει ένα σύνολο σημείων από ένα σύστημα συντεταγμένων σε ένα άλλο. Όταν είναι γνωστές αντιστοιχίες σημείων $(\mathbf{p}_i, \mathbf{m}_i)$, $i = 1, \dots, n$, το πρόβλημα είναι γνωστό ως πρόβλημα του \emph{απόλυτου προσανατολισμού}: ζητούνται στροφή $R$, μετάθεση $\mathbf{t}$ και, προαιρετικά, κλίμακα $s$ που ελαχιστοποιούν
\begin{equation}
  E(R, \mathbf{t}, s) = \sum_{i=1}^{n} \lVert s R \mathbf{p}_i + \mathbf{t} - \mathbf{m}_i \rVert^2 .
  \label{eq:absolute-orientation}
\end{equation}
Το πρόβλημα έχει κλειστές λύσεις. Ο \en{Horn} τη διατύπωσε με μοναδιαία τετραδόνια \cite{horn1987closed}. Οι \en{Arun} κ.ά. και ο \en{Kabsch} μέσω της ανάλυσης ιδιαζουσών τιμών (SVD) του πίνακα συνδιακύμανσης των κεντραρισμένων σημείων \cite{arun1987leastsquares,kabsch1976solution}. Ο \en{Umeyama} έδωσε τη γενική λύση με κλίμακα και διόρθωσε την περίπτωση όπου η SVD επιστρέφει κατοπτρισμό αντί για στροφή \cite{umeyama1991leastsquares}. Σε όλες τις περιπτώσεις η μετάθεση προκύπτει από τα κέντρα βάρους: $\mathbf{t} = \bar{\mathbf{m}} - sR\bar{\mathbf{p}}$.

Η ποιότητα μιας ευθυγράμμισης αποτιμάται με τα \emph{υπόλοιπα} $\mathbf{r}_i = sR\mathbf{p}_i + \mathbf{t} - \mathbf{m}_i$ και τη ρίζα του μέσου τετραγώνου τους (RMS). Το RMS στα ίδια τα σημεία ευθυγράμμισης είναι όμως αισιόδοξο, γιατί η λύση έχει προσαρμοστεί σε αυτά. Η πραγματική ακρίβεια μετριέται σε \emph{ανεξάρτητα σημεία ελέγχου} (\en{check points}), που δεν συμμετέχουν στη λύση, όπως γίνεται στην τοπογραφία. Αυτή η διάκριση ακολουθείται στο πείραμα της \cref{sec:exp-alignment}.

Όταν δεν υπάρχουν γνωστές αντιστοιχίες αλλά δύο νέφη σημείων ή επιφάνειες, χρησιμοποιείται ο επαναληπτικός αλγόριθμος ICP \cite{besl1992icp}, που εναλλάσσει την εύρεση πλησιέστερων σημείων με την κλειστή λύση της \cref{eq:absolute-orientation}. Η εργασία δεν τον χρησιμοποιεί, αλλά τον αναφέρει ως μελλοντική επέκταση (\cref{sec:future}).

\section{Παρεμβολή χωρικών δεδομένων}
\label{sec:bg-interpolation}

Οι αισθητήρες μετρούν σε διακριτά σημεία. Για να απεικονιστεί ένα μέγεθος, π.χ. η μετακίνηση, πάνω σε μια συνεχή επιφάνεια, πρέπει να εκτιμηθεί η τιμή του ανάμεσα στα σημεία μέτρησης.

Η μέθοδος \emph{αντίστροφης στάθμισης απόστασης} (IDW) του \en{Shepard} εκτιμά την τιμή σε ένα σημείο $\mathbf{x}$ ως σταθμισμένο μέσο όρο όλων των μετρήσεων $z_i$ \cite{shepard1968interpolation}:
\begin{equation}
  \hat{z}(\mathbf{x}) = \frac{\sum_i w_i(\mathbf{x})\, z_i}{\sum_i w_i(\mathbf{x})}, \qquad
  w_i(\mathbf{x}) = \lVert \mathbf{x} - \mathbf{x}_i \rVert^{-p},
  \label{eq:idw}
\end{equation}
με εκθέτη $p$ συνήθως ίσο με 2. Είναι απλή και ορίζεται παντού, αλλά δημιουργεί χαρακτηριστικά «μάτια» γύρω από κάθε μέτρηση, γιατί η επιφάνεια επιπεδώνεται κοντά σε κάθε σημείο, και επηρεάζεται από μακρινά σημεία.

Η \emph{τριγωνοποίηση Delaunay} ενός συνόλου σημείων στο επίπεδο είναι η τριγωνοποίηση όπου κανένα σημείο δεν βρίσκεται μέσα στον περιγεγραμμένο κύκλο κάποιου τριγώνου. Μεγιστοποιεί την ελάχιστη γωνία και αποφεύγει τα στενόμακρα τρίγωνα \cite{deberg2008computational}. Πάνω σε αυτήν, η \emph{γραμμική παρεμβολή} δίνει σε κάθε σημείο τον βαρυκεντρικό συνδυασμό των τιμών των τριών κορυφών του τριγώνου που το περιέχει. Το αποτέλεσμα είναι συνεχές, περνά ακριβώς από τις μετρήσεις και χρησιμοποιεί μόνο τους γειτονικούς αισθητήρες. Μειονέκτημα είναι ότι δεν ορίζεται έξω από το κυρτό περίβλημα των σημείων. Πιο σύνθετες γεωστατιστικές μέθοδοι, όπως το \en{kriging}, μοντελοποιούν τη χωρική συσχέτιση αλλά απαιτούν περισσότερα δεδομένα για την εκτίμηση του βαριογράμματος.

\section{Απόδοση γραφικών σε αυτόνομες συσκευές}
\label{sec:bg-rendering}

Μια εφαρμογή εικονικής ή μικτής πραγματικότητας πρέπει να διατηρεί σταθερό ρυθμό καρέ ίσο με τον ρυθμό ανανέωσης της οθόνης. Αλλιώς η εικόνα τρεμοπαίζει και ο χρήστης μπορεί να νιώσει δυσφορία. Στα 72\,Hz κάθε καρέ έχει προϋπολογισμό \SI{13.9}{ms}, και στα 90\,Hz \SI{11.1}{ms}, για δύο εικόνες (μία ανά μάτι).

Σε μια κινητή GPU το κόστος ενός καρέ καθορίζεται κυρίως από τρεις παράγοντες. Ο πρώτος είναι το \emph{πλήθος των draw calls}, δηλαδή των ξεχωριστών εντολών σχεδίασης που στέλνει η CPU. Μειώνεται με την ομαδοποίηση (\en{batching}) αντικειμένων που μοιράζονται υλικό, και με τη σχεδίαση και των δύο ματιών σε ένα πέρασμα (\en{single-pass instanced}). Ο δεύτερος είναι το \emph{πλήθος των τριγώνων και των κορυφών}. Ο τρίτος είναι το \emph{κόστος ανά pixel}, που αυξάνεται απότομα με την ημιδιαφάνεια, γιατί κάθε ημιδιαφανής επιφάνεια ζωγραφίζεται πάνω από τις προηγούμενες (\en{overdraw}). Το \emph{foveated rendering} μειώνει την ανάλυση στην περιφέρεια του οπτικού πεδίου, όπου το ανθρώπινο μάτι διακρίνει λιγότερη λεπτομέρεια, και ελευθερώνει χρόνο GPU \cite{patney2016foveated}.

Τα μοντέλα φωτογραμμετρίας έχουν συχνά εκατομμύρια τρίγωνα και υφές υψηλής ανάλυσης, πολύ πάνω από όσα αντέχει μια αυτόνομη συσκευή. Η \emph{απλοποίηση πλέγματος} μειώνει τα τρίγωνα με τη μικρότερη δυνατή γεωμετρική αλλοίωση. Η μέθοδος των \en{Garland} και \en{Heckbert} συμπτύσσει διαδοχικά ακμές, επιλέγοντας κάθε φορά αυτήν με το μικρότερο τετραγωνικό σφάλμα (QEM) ως προς τα επίπεδα των γειτονικών τριγώνων \cite{garland1997qem}. Η απόκλιση του απλοποιημένου πλέγματος από το αρχικό μετριέται συνήθως με την \emph{απόσταση Hausdorff}, όπως στο εργαλείο \en{Metro} \cite{cignoni1998metro}. Και τα δύο είναι διαθέσιμα στο MeshLab \cite{cignoni2008meshlab}.

\section{Παρακολούθηση κατασκευών και πλατφόρμες IoT}
\label{sec:bg-shm}

Η παρακολούθηση δομικής ακεραιότητας είναι η διαδικασία εντοπισμού μεταβολών σε μια κατασκευή μέσω συστηματικών μετρήσεων στον χρόνο \cite{farrar2007shm}. Σε γεωτεχνικά έργα και υποδομές χρησιμοποιούνται, μεταξύ άλλων:
\begin{itemize}
  \item \textbf{τοπογραφικοί στόχοι} (πρίσματα), που μετρώνται με γεωδαιτικό σταθμό και δίνουν θέση $(x, y, z)$ και μετακινήσεις $(dx, dy, dz)$ ως προς μια μηδενική μέτρηση·
  \item \textbf{κλισιόμετρα}, για τη μεταβολή κλίσης·
  \item \textbf{ρωγμόμετρα}, για το άνοιγμα ρωγμών και αρμών·
  \item \textbf{πιεζόμετρα}, για τη στάθμη ή την πίεση του νερού.
\end{itemize}
Σε κάθε αισθητήρα ορίζονται \emph{όρια} σε βαθμίδες (π.χ. αναφορά, προειδοποίηση, συναγερμός). Όταν μια μέτρηση ξεπερνά ένα όριο, η πλατφόρμα δημιουργεί συμβάν συναγερμού.

Οι μετρήσεις συλλέγονται σε μια πλατφόρμα παρακολούθησης, που στην παρούσα εργασία είναι το Terraweb. Οργανώνει τα δεδομένα σε έργα, περιοχές, αισθητήρες, χρονοσειρές μετρήσεων, όρια και ενεργά συμβάντα, και τα εκθέτει μέσω διεπαφής REST \cite{fielding2000rest} με αυθεντικοποίηση JWT. Η πρόσβαση της συσκευής σε αυτή τη διεπαφή περιγράφεται στην \cref{sec:comm-terraweb}.

\review{Επιβεβαίωσε την περιγραφή του Terraweb (οργάνωση σε έργα/περιοχές, όρια σε βαθμίδες, ορολογία των βαθμίδων). Αν η πλατφόρμα έχει δημόσια τεκμηρίωση ή είναι προϊόν εταιρείας, δες με τον επιβλέποντα \textbf{πώς επιτρέπεται να αναφερθεί} (όνομα, εταιρεία, εμπιστευτικότητα).}
